\begin{theorem}
  \label{ent:thm:fermat}
  If $p$ is prime and $p\nmid a$, then $a^{p-1}\equiv1\pmod p$.

  \begin{lean}
    theorem fermat (p : ℕ) (hp : Chapter02.isPrime p) (a : ℤ) (ha : ¬(p : ℤ) ∣ a) :
        a ^ (p - 1) ≡ 1 [ZMOD (p : ℤ)]
  \end{lean}
\end{theorem}

\begin{proof}
  Multiplication by the nonzero class of $a$ permutes the nonzero residue classes.
  Multiplying all these equalities yields $a^{p-1}P=P$, where $P$ is the product of all nonzero
  classes.
  Each factor is invertible, so cancellation gives the result.
  The auxiliary permutation-product lemma implements this finite argument directly.

  \begin{lean}
    let : Fact p.Prime := ⟨Nat.prime_def.mpr hp⟩

    have hne : (a : ZMod p) ≠ 0 := by
      intro h
      exact ha ((ZMod.intCast_zmod_eq_zero_iff_dvd a p).1 h)

    have h := pow_card_by_permutation (Units.mk0 (a : ZMod p) hne)

    have hv := congrArg (fun u : (ZMod p)ˣ => (u : ZMod p)) h
    simp only [Units.val_pow_eq_pow_val, Units.val_mk0, Units.val_one,
      Fintype.card_units, ZMod.card] at hv

    apply (ZMod.intCast_eq_intCast_iff _ _ p).1
    simpa only [Int.cast_pow, Int.cast_one] using hv
  \end{lean}
\end{proof}

\begin{theorem}
  \label{ent:thm:fermat-all}
  For every prime $p$ and integer $a$, one has $a^p\equiv a\pmod p$.

  \begin{lean}
    theorem fermat_all (p : ℕ) (hp : Chapter02.isPrime p) (a : ℤ) :
        a ^ p ≡ a [ZMOD (p : ℤ)]
  \end{lean}
\end{theorem}

\begin{proof}
  If $p\mid a$, both sides are zero modulo $p$.
  Otherwise multiply Fermat\textquotesingle s congruence by $a$.

  \begin{lean}
    have hpos : 0 < p := by
      have := hp.1
      omega
    by_cases hpa : (p : ℤ) ∣ a

    · have hz := Chapter03.modEq_pow ((Chapter03.modEq_zero_iff _ _).2 hpa) p
      rw [zero_pow (ne_of_gt hpos)] at hz
      exact Chapter03.modEq_trans hz
        (Chapter03.modEq_symm ((Chapter03.modEq_zero_iff _ _).2 hpa))

    · have h := Chapter03.modEq_mul (fermat p hp a hpa) (Chapter03.modEq_refl (p : ℤ) a)
      simpa only [← pow_succ, Nat.sub_add_cancel hpos, one_mul] using h
  \end{lean}
\end{proof}

\begin{theorem}
  \label{ent:thm:composite-of-not-fermat}
  If $n>1$ and $a^n\not\equiv a\pmod n$ for some integer $a$, then $n$ is composite.

  \begin{lean}
    theorem composite_of_not_fermat (n : ℕ) (hn : 1 < n) (a : ℤ)
        (h : ¬a ^ n ≡ a [ZMOD (n : ℤ)]) : Chapter02.isComposite n
  \end{lean}
\end{theorem}

\begin{proof}
  If $n$ were prime, the unrestricted Fermat congruence would contradict the hypothesis.

  \begin{lean}
    ⟨hn, fun hp => h (fermat_all n hp a)⟩
  \end{lean}
\end{proof}

\begin{theorem}
  \label{ent:thm:fermat-two-primes}
  Let $p\ne q$ be primes.
  If $a^p\equiv a\pmod q$ and $a^q\equiv a\pmod p$, then $a^{pq}\equiv a\pmod{pq}$.

  \begin{lean}
    theorem fermat_two_primes (p q : ℕ) (hp : Chapter02.isPrime p)
        (hq : Chapter02.isPrime q) (hne : p ≠ q) (a : ℤ)
        (hpq : a ^ p ≡ a [ZMOD (q : ℤ)]) (hqp : a ^ q ≡ a [ZMOD (p : ℤ)]) :
        a ^ (p * q) ≡ a [ZMOD ((p * q : ℕ) : ℤ)]
  \end{lean}
\end{theorem}

\begin{proof}
  Raise the congruence modulo $p$ to the $p$th power, and the one modulo $q$ to the $q$th power.
  Fermat\textquotesingle s theorem reduces their right-hand sides to $a$.
  The two distinct primes are coprime, so their product divides $a^{pq}-a$.

  \begin{lean}
    have h₁ := Chapter03.modEq_trans (Chapter03.modEq_pow hqp p) (fermat_all p hp a)

    have h₂ := Chapter03.modEq_trans (Chapter03.modEq_pow hpq q) (fermat_all q hq a)
    rw [← pow_mul, Nat.mul_comm q p] at h₁
    rw [← pow_mul] at h₂
    apply (Chapter03.modEq_iff_dvd_sub _ _ _).2
    rw [Nat.cast_mul]
    exact Chapter01.mul_dvd_of_coprime p q _ (gcd_distinct_primes p q hp hq hne)
      ((Chapter03.modEq_iff_dvd_sub _ _ _).1 h₁) ((Chapter03.modEq_iff_dvd_sub _ _ _).1 h₂)
  \end{lean}
\end{proof}

\begin{definition}
  A composite natural number $n$ is a \emph{pseudoprime to base $a$} if $a^n\equiv a\pmod n$.
  The term pseudoprime without a specified base refers to base two.

  \begin{lean}
    def IsPseudoprime (a : ℤ) (n : ℕ) : Prop :=
      Chapter02.isComposite n ∧ a ^ n ≡ a [ZMOD (n : ℤ)]
  \end{lean}
\end{definition}

\begin{definition}
  An \emph{absolute pseudoprime}, or Carmichael number, is a composite $n$ satisfying $a^n\equiv
    a\pmod n$ for every integer $a$.

  \begin{lean}
    def IsAbsolutePseudoprime (n : ℕ) : Prop :=
      Chapter02.isComposite n ∧ ∀ a : ℤ, a ^ n ≡ a [ZMOD (n : ℤ)]
  \end{lean}
\end{definition}

\begin{definition}
  A positive integer is \emph{square-free} when no square of a prime divides it.
  In particular, one is square-free.

  \begin{lean}
    def IsSquareFree (n : ℕ) : Prop :=
      0 < n ∧ ∀ p : ℕ, Chapter02.isPrime p → ¬p ^ 2 ∣ n
  \end{lean}
\end{definition}

\begin{theorem}
  \label{ent:thm:pseudoprime-two-iff}
  For an odd composite $n$, the base-two pseudoprime condition is equivalent to
  $2^{n-1}\equiv1\pmod n$.

  \begin{lean}
    theorem pseudoprime_two_iff (n : ℕ) (hn : Odd n) (hc : Chapter02.isComposite n) :
        IsPseudoprime 2 n ↔ (2 : ℤ) ^ (n - 1) ≡ 1 [ZMOD (n : ℤ)]
  \end{lean}
\end{theorem}

\begin{proof}
  An odd integer is coprime to two: writing $n=2k+1$ gives the Bezout identity $1=n-2k$.
  Cancel two in one direction and multiply by two in the other.

  \begin{lean}
    have hpos : 0 < n := by
      have := hc.1
      omega

    constructor

    · intro h
      apply Chapter03.modEq_cancel (gcd_odd_two n hn)
      simpa only [← pow_succ', Nat.sub_add_cancel hpos, mul_one] using h.2

    · intro h
      refine ⟨hc, ?_⟩

      have hh := Chapter03.modEq_mul h (Chapter03.modEq_refl (n : ℤ) 2)
      simpa only [← pow_succ, Nat.sub_add_cancel hpos, one_mul] using hh
  \end{lean}
\end{proof}

\begin{definition}
  The Mersenne number with exponent $n$ is $M_n=2^n-1$.

  \begin{lean}
    def mersenne (n : ℕ) : ℕ := 2 ^ n - 1
  \end{lean}
\end{definition}

\begin{theorem}
  \label{ent:thm:mersenne-pseudoprime}
  If $n$ is a base-two pseudoprime, then $M_n$ is a larger base-two pseudoprime.
  In particular, this holds for every odd pseudoprime.

  \begin{lean}
    theorem mersenne_pseudoprime (n : ℕ) (h : IsPseudoprime 2 n) :
        IsPseudoprime 2 (mersenne n) ∧ n < mersenne n
  \end{lean}
\end{theorem}

\begin{proof}
  Since $n$ is composite, it has a proper divisor $d>1$.
  The identity $d\mid n$ gives $2^d-1\mid2^n-1$, and this is a nontrivial proper divisor.
  Thus $M_n$ is composite.
  Induction also gives $2^n-1>n$ for $n\ge2$.

  \begin{lean}
    have hn : 2 ≤ n := h.1.1

    have hlarge := lt_mersenne n hn

    have hpow : 2 ≤ 2 ^ n := by
      have hh : 1 ≤ 2 ^ (n - 1) := by
        have := pow_pos (by decide : 0 < (2 : ℕ)) (n - 1)
        omega

      have he : 2 ^ n = 2 ^ (n - 1) * 2 := by
        rw [← pow_succ, Nat.sub_add_cancel (by omega : 1 ≤ n)]
      omega
  \end{lean}

  The pseudoprime congruence gives $n\mid2^n-2=M_n-1$.
  Write $M_n-1=nk$.
  Since $2^n\equiv1\pmod{M_n}$, raising to the $k$th power and then multiplying by two proves
  $2^{M_n}\equiv2\pmod{M_n}$.

  \begin{lean}
    have hdiv : n ∣ mersenne n - 1 := by
      apply Int.natCast_dvd_natCast.mp

      have hh := (Chapter03.modEq_iff_dvd_sub _ _ _).1 h.2

      have he : ((mersenne n - 1 : ℕ) : ℤ) = (2 : ℤ) ^ n - 2 := by
        rw [mersenne, Nat.sub_sub, Nat.cast_sub hpow, Nat.cast_pow, Nat.cast_ofNat]
      rwa [he]

    obtain ⟨k, hk⟩ := hdiv

    have hbase : (2 : ℤ) ^ n ≡ 1 [ZMOD ((mersenne n : ℕ) : ℤ)] := by
      apply (Chapter03.modEq_iff_dvd_sub _ _ _).2
      refine ⟨1, ?_⟩
      simp only [mersenne, Nat.cast_sub (by omega : 1 ≤ 2 ^ n),
        Nat.cast_pow, Nat.cast_ofNat, Nat.cast_one, mul_one]

    have hpower := Chapter03.modEq_pow hbase k
    rw [← pow_mul, ← hk, one_pow] at hpower

    have hfinal := Chapter03.modEq_mul hpower
      (Chapter03.modEq_refl ((mersenne n : ℕ) : ℤ) 2)
    rw [← pow_succ, Nat.sub_add_cancel (by omega : 1 ≤ mersenne n), one_mul] at hfinal

    exact ⟨⟨mersenne_composite n h.1, hfinal⟩, hlarge⟩
  \end{lean}
\end{proof}

\begin{theorem}
  \label{ent:thm:infinite-odd-pseudoprimes}
  There are infinitely many odd base-two pseudoprimes.

  \begin{lean}
    theorem infinite_odd_pseudoprimes : {n : ℕ | Odd n ∧ IsPseudoprime 2 n}.Infinite
  \end{lean}
\end{theorem}

\begin{proof}
  Start from $341$ and repeatedly apply the Mersenne construction.
  Every new value is odd and strictly larger, so an odd pseudoprime exists above every natural
  bound.
  An unbounded subset of the natural numbers is infinite.

  \begin{lean}
    apply Set.infinite_of_forall_exists_gt
    intro B

    obtain ⟨n, hB, hodd, hn⟩ := exists_odd_pseudoprime_gt B

    exact ⟨n, ⟨hodd, hn⟩, hB⟩
  \end{lean}
\end{proof}

\begin{theorem}
  \label{ent:thm:absolutepseudoprime-of-prime-product}
  Let $n$ be a composite product of distinct primes.
  If $p-1\mid n-1$ for every prime in the product, then $n$ is an absolute pseudoprime.

  \begin{lean}
    theorem absolutePseudoprime_of_prime_product (s : Finset ℕ)
        (hp : ∀ p ∈ s, Chapter02.isPrime p)
        (hc : Chapter02.isComposite (∏ p ∈ s, p))
        (hdiv : ∀ p ∈ s, p - 1 ∣ (∏ q ∈ s, q) - 1) :
        IsAbsolutePseudoprime (∏ p ∈ s, p)
  \end{lean}
\end{theorem}

\begin{proof}
  Fix $a$.
  For each prime factor $p$, either $p\mid a$, or Fermat\textquotesingle s theorem and $p-1\mid
    n-1$ give $a^{n-1}\equiv1\pmod p$.
  In either case $a^n\equiv a\pmod p$.
  Pairwise coprimality combines these congruences modulo their product.

  \begin{lean}
    refine ⟨hc, ?_⟩
    intro a
    apply (Chapter03.modEq_iff_dvd_sub _ _ _).2
    rw [Nat.cast_prod]
    apply Chapter03.prod_dvd_of_pairwise s (fun p => (p : ℤ))

    · intro p hps

      have := (hp p hps).1
      omega

    · intro p hps q hqs hne
      exact gcd_distinct_primes p q (hp p hps) (hp q hqs) hne

    · intro p hps
      exact (Chapter03.modEq_iff_dvd_sub _ _ _).1
        (fermat_of_sub_one_dvd p _ (hp p hps) (by
            have := hc.1
            omega) (hdiv p hps) a)
  \end{lean}
\end{proof}

\begin{theorem}
  \label{ent:thm:absolutepseudoprime-squarefree}
  Every absolute pseudoprime is square-free.

  \begin{lean}
    theorem absolutePseudoprime_squareFree (n : ℕ) (h : IsAbsolutePseudoprime n) :
        IsSquareFree n
  \end{lean}
\end{theorem}

\begin{proof}
  If $p^2\mid n$, use the defining congruence with base $p$.
  Since $n\ge2$, both $p^n$ and $n$ are divisible by $p^2$, so $p^2\mid p$.
  Cancelling $p$ would make a prime divide one, a contradiction.

  \begin{lean}
    refine ⟨by
        have := h.1.1
        omega, ?_⟩
    intro p hp hpn

    have hn : 2 ≤ n := h.1.1

    obtain ⟨k, hk⟩ := (Chapter03.modEq_iff_dvd_sub _ _ _).1 (h.2 p)
    obtain ⟨t, ht⟩ := hpn

    have ht' : (n : ℤ) = (p : ℤ) ^ 2 * t := by exact_mod_cast ht

    have hpow : (p : ℤ) ^ n = (p : ℤ) ^ 2 * (p : ℤ) ^ (n - 2) := by
      rw [← pow_add, Nat.add_sub_cancel' hn]

    have hcancel : (p : ℤ) * ((p : ℤ) ^ (n - 2) - (t : ℤ) * k) = 1 := by
      have hpne : (p : ℤ) ≠ 0 := by
        have := hp.1
        omega

      apply mul_left_cancel₀ hpne
      rw [ht', hpow] at hk
      nlinarith [hk]

    have hd : p ∣ 1 := Int.natCast_dvd_natCast.mp
      ⟨(p : ℤ) ^ (n - 2) - (t : ℤ) * k, hcancel.symm⟩

    obtain ⟨u, hu⟩ := hd

    have := hp.1

    have hu' : 0 < u := by nlinarith
    nlinarith
  \end{lean}
\end{proof}
